\subsection{Dostrajanie modeli (fine-tuning) i metody parameter-efficient}
\label{subsec:llm_acr_finetuning}

W podejściu opartym na dostrajaniu (\emph{fine-tuning}) przyjmuje się, że część ograniczeń promptowania 
(zwłaszcza to, że jakość bywa nierówna i czasami ciężko zapewnić oczekiwanej formy komentarza) da się ograniczyć
przez dopasowanie modelu do rozkładu danych pochodzących z procesu przeglądu kodu. Z punktu widzenia procesów
wytwarzania oprogramowania, wpasowuje się to w szerszą obserwację, że LLM w tym obszarze wykorzystuje się w wielu
zadaniach, lecz wykorzystanie bardzo ogólnego modelu do wyspecjalizowanego procesu, jakim jest code review,
zwykle wymaga dopasowania jednocześnie danych i celu uczenia \cite{Zhang2024SurveyLLMSE}.

\paragraph{Parameter-efficient fine-tuning jako kompromis wdrożeniowy.}
Dostrajanie oszczędne parametrowo (PEFT) jest jednym z ważniejszych kierunków w tym obszarze. Jest to podejście
w którym model adaptuje się do zadania bez modyfikowania pełnego zestawu wag. Lu i in. dobrze obrazują to podejście
poprzez LLaMA-Reviewer i pokazują, że zamiast trzymać osobne, ciężkie kopie modelu można przechowywać po prostu elementy dostrojenia.
Dla LLaMA 6.7B w ich opisie LoRA to ok. \(\sim 8.4\) mln trenowalnych parametrów i \(\sim 16\) MB, 
a prefix-tuning schodzi do ok. \(\sim 1.2\) mln parametrów i \(\sim 2.4\) MB \cite{Lu2023LLaMAReviewer}. PEFT nie jest
tam użyte tylko do jednego kroku, ale do kilku zadań w pipeline review (m.in. predykcji potrzeby review, 
generowania komentarza oraz \emph{code refinement}) \cite{Lu2023LLaMAReviewer}.
W tym miejscu istotne wydaje się takie wieloetapowie wykorzystanie tego podejścia, bo recenzja kodu w realnym narzędziu
zwykle nie kończy się na samym wygenerowaniu tekstu, tylko dochodzą kolejne etapy i decyzje. Podobny wniosek pojawia się
też w badaniach na zbiorze CodeReviewer: Haider i in. raportują, że PEFT (konkretnie QLoRA) daje mierzalne zyski w RCG, np. BLEU-4 rośnie do 5.58 
dla Code Llama 7B (vs 4.28 dla baseline CodeReviewer 223M), a BERTScore 
do 0.8483 dla Llama 3.1 8B (vs 0.8348) \cite{Haider2024PromptingFineTuningRCG}. Patrząc na to 
z perspektywy zdefiniowania state-of-the-art, wygląda na to, że fine-tuning nie jest "zarezerwowany" wyłącznie 
dla dużych laboratoriów i centr obliczeniowych, chociaż nadal łatwo tu o rozjazd wyników, jeśli dane 
albo protokół oceny są źle dobrane \cite{Tufano2024CodeReviewAutomation}.

\paragraph{Fine-tuning z naciskiem na zrozumiałość i strukturalną kompletność review.}
Faktem jest, że ocena jakości automatycznego code review nie powinna kończyć się na pytaniu "czy komentarz leksykalnie pasuje do referencji".
W praktyce w review liczy się raczej to, czy da się z komentarza wyciągnąć: gdzie jest problem, z czego on wynika i co można z 
tym zrobić (czyli jaka jest propozycja naprawy). Yu i in. projektują Carllm właśnie pod kątem poziomu jasności i
zrozumienia komentarzy i budują dane w oparciu o szerszy kontekst dyskusji review. Widać tu dość konkretne decyzje inżynierskie: z 1\,793 projektów 
open-source wybrano te spełniające kryteria aktywności (m.in. >1000 PR) oraz licencji (MIT/Apache 
2.0), a następnie z 15\,000 rekordów wyprowadzono 10\,554 (77.02\%) przykładów wysokiej jakości; 
po dalszym czyszczeniu zostało 9\,728 rekordów. Dalej dodano przykłady negatywne w proporcji 
1:1, co dało 19\,456 instancji treningowych \cite{Yu2024Carllm}. Autorzy zwracają też uwagę, że 
LoRA może obejmować bardzo mały ułamek parametrów (np. \(r=8\) to ok. 0.05\% 
parametrów modelu), więc sama adaptacja jest wykonalna nawet przy ograniczonych zasobach \cite{Yu2024Carllm}.
Wyniki benchmarkingowe pokazują istnienie ryzyka regresji (pogorszenia wynikowego kodu) 
nawet przy silnych modelach, więc nacisk na czytelność i sensowne uzasadnienie ma 
znaczenie nie tylko "użytkowe", ale też związane z bezpieczeństwem procesu \cite{Cihan2025EvaluatingLLMsForCodeReview}.

\paragraph{Jakość danych: destylacja "pożądanych" komentarzy jako warstwa kontrolna.}
Poprawa jakości samego sygnału uczącego jest dodatkowo ważnym wątkiem uzupełniającym. Yu i in. zwracają uwagę, że
komentarze review nie mają tej samej wartości, jeśli model ma się uczyć generowania uwag faktycznie prowadzących do
naprawy kodu. Autorzy wprowadzają pojęcie \emph{Desired Review Comments} (DRC) i metodę Desiview która przypisuje komentarzom
wynik pożądalności (\emph{desiredness score}). W praktyce jest to różnica perplexity (wariant z komentarzem vs bez komentarza), 
a wynik jest dodatkowo stabilizowany przez głosowanie czterech modeli (m.in. CodeLlama-13B, Meta-Llama-3-8B) 
\cite{Yu2025Desiview}. Na ręcznie oznaczanej próbce 600 rekordów Desiview osiąga 86.67\% accuracy oraz 
84.44\% F1 i wypada lepiej od GPT-4o (76.50\% accuracy, 71.05\% F1) \cite{Yu2025Desiview}. Autorzy pokazują też na przykładzie CodeReviewer,
że DRC to tylko ok. 43\% danych (np. 64\,934 z 150\,406 w treningu), więc "surowy" zbiór 
jest w dużej części zaszumiony, patrząc na cel generowania komentarzy prowadzących do 
realnych poprawek \cite{Yu2025Desiview}. Autorzy dostroili następnie model na wyczyszczonych danych, co poprawiło
BLEU-4 (np. dla LLaMA-3: 8.33 \(\rightarrow\) 11.87 po fine-tuningu, a po dodatkowym dopasowaniu do 13.13). Zwiększyło
też to odsetek ocen bardzo bliskich człowiekowi \cite{Yu2025Desiview}. Z perspektywy projektu wniosek jest prosty: samo dokładanie danych nie wystarczy,
a fine-tuning warto oprzeć na selekcji/filtracji przykładów, bo inaczej łatwo przepuścić do uczenia komentarze, które są mało 
pomocne.

\paragraph{Architektury wielozadaniowe i wykorzystanie MoE w procesie review.}
Tang i in. zwracają uwagę na problem: uczenie poszczególnych elementów pipeline'u 
recenzji kodu w izolacji sprawia, że wiedza nie przepływa między zadaniami tak, 
jak powinna. Chodzi głównie o relację między tym, czy system w ogóle uzna, że warto coś recenzować, a tym, jak potem faktycznie generuje 
dany komentarz. Wewnątrz mamy ekspertów wyspecjalizowanych w trzech konkretnych zadaniach: przewidywaniu konieczności recenzji (necessity 
prediction), generowaniu komentarzy oraz poprawianiu kodu (refinement). Całość spina mechanizm routowania Top-2\footnote{każdego wejścia wybierane są dwa najbardziej odpowiednie eksperty - modele, a ich wyniki są łączone}
wraz z load-balancingiem, co ma zapewniać odpowiednią efektywność \cite{Tang2025MoEReviewer}. 

Podejście to faktycznie ma dobre zastosowanie, co widać na wynikach na zbiorze CodeReviewer. Dla predykcji konieczności recenzji uzyskano 74,3\% accuracy (F1 na poziomie 
73,2\%). Przy samym generowaniu komentarzy BLEU wyniosło 11,62, a dla zadania refinement 
wynik ten dochodzi do 87,37 \cite{Tang2025MoEReviewer}. Bardzo ważnym aspektem tej pracy, zwłaszcza 
patrząc na to od strony czysto implementacyjnej czy bizneowej, jest wydajność obliczeniowa. Dzięki strukturze 
MoE model aktywuje tylko około 2 mld parametrów na token, co daje 
koszt rzędu 16,38 TFLOPs. To duży przeskok, bo wspomniany w pracy LLaMA-Reviewer 
(7B) wymaga aż 57,34 TFLOPs \cite{Tang2025MoEReviewer}. Ma to kluczowe znaczenie, bo narzędzia 
typu automated code review muszą przecież działać w ramach konkretnych budżetów czasowych i pieniężnych w procesach 
CI/CD. W takich warunkach opóźnienie inferencji i ogólny koszt są w zasadzie 
tak samo ważne, jak same metryki jakościowe modelu.

\paragraph{Wnioski z przeglądu literatury: fine-tuning a realia procesu review.}
W rzeczywistych warunkach code review jest procesem bardzo złożonym i wielowymiarowym. Sprawia to, że
ciężko jest sprowadzić to zadanie do prostego modelu 1:1, czyli relacji typu $diff \rightarrow$ pojedynczy komentarz.
Dobrze pokazują to Li i in. w pracy nad CodeDoctor - z ich analizy zbioru Turzo2023 wynika, że aż 78,7\% 
komentarzy dotyczyło co najmniej dwóch różnych kategorii. W przypadku zbioru Tufano2021 było 
to 56,71\%. Średnio na jeden komentarz przypadały dwie kategorie, a w ponad 
10\% przypadków było ich cztery lub nawet więcej \cite{Li2025CodeDoctor}. Ci sami autorzy 
zauważają, że w Turzo2023 zdecydowana większość procesów review (92,7\%) zajmowała programistom ponad 
pół godziny, co w zasadzie jest najlepszym argumentem za tym, żeby szukać 
metod automatyzacji wspierających recenzenta \cite{Li2025CodeDoctor}.

W pracach zbierających aktualny stan wiedzy, zauważyć można sporą dawkę sceptycyzmu wobec stosowania samych metryk.
Tufano i in. \cite{Tufano2024CodeReviewAutomation} trafnie wskazują, 
że bez dokładnej analizy jakościowej i kontroli danych łatwo wyciągnąć błędne wnioski. Bazowanie na samych metrykach
może ukrywać problemy z mapowaniem komentarzy na zmiany w kodzie. Może też maskować fakt, że dany model radzi sobie
bardzo dobrze tylko w tych najprostszych, powtarzalnych i często występujących przypadkach.

Bardzo ważnym aspektem jest także kwestia bezpieczeństwa oprogramowania. Germano i in. przeprowadzili systematyczny przegląd ponad dwustu publikacji z okresu 2018-sierpień 
2025 i widać tam ogromną dysproporcję w podejmowanych tematach. Aż 91,3\% prac 
skupia się wyłącznie na wykrywaniu podatności Z kolei naprawianie błędów
czy ich wyjaśnianie to odpowiednio tylko 11,1\% i 5,3\% publikacji \cite{Germano2025VulnSLR}. 
To pokazuje lukę w badaniach - w systemach automatycznego przeglądu kodu nie wystarczy przecież 
samo "wskazanie palcem" problemu. Żeby narzędzie było użyteczne, musi oferować zrozumiałe uzasadnienie 
i konkretne propozycje zmian, co zresztą potwierdzają autorzy prac nad podejściami DRC 
i comprehensibility \cite{Yu2024Carllm,Yu2025Desiview}.

Wdrożenie w przemyśle jest równie istotnym aspektem. Ciekawym przykładem są doświadczenia z firmy Ericsson, opisane przez 
Ramesha \cite{Ramesh2025EricssonExperience}. Tam zamiast kosztownego fine-tuningu postawiono na "lekkie" podejście zintegrowane bezpośrednio 
z workflow (Gerrit, VS Code). Wykorzystali oni parser Tree-sitter do wyciągania kontekstu 
z AST (tzw. enclosing method), co pokazuje, że w praktyce wybór modelu to zawsze kompromis. 
Wynika z tego, że w praktyce należy wyważyć jakość wyników, koszty utrzymania 
i to, jak trudno będzie zintegrować narzędzie z istniejącym procesem wytwórczym \cite{Ramesh2025EricssonExperience}.

\paragraph{Podsumowanie.}
Dostrajanie modeli (szczególnie PEFT) jest dziś jednym z ważniejszych elementów state-of-the-art w 
automatycznym \emph{code review}. Pozwala utrzymać bardziej przewidywalny styl i dość 
spójną strukturę komentarzy. Pozwala też lepiej dopasować zachowanie modelu do tego, jak 
wyglądają dane w konkretnym repozytorium. Ekuteczność takiego podejścia mocno zależy od: jego jakości, oraz udziale komentarzy, które 
faktycznie uznajemy za pożądane. Do tego dochodzi kwestia tego, jak w ogóle 
mierzymy jakość (protokół ewaluacji) oraz koszty wdrożenia, które czasem są niesłusznie pomijane.

Część typowych problemów w \emph{review} (braki kontekstu projektowego, halucynacje albo zgodność z regułami jakości) bywa
łatwiej ograniczać metodami hybrydowymi i/lub przez retrieval, zamiast samym dostrajaniem. W literaturze pojawia się to m.in. jako klasyczny 
RAG dopasowany do \emph{review} \cite{Meng2025RARe}, ale też jako integracja z analizą statyczną 
(RAG oparty o wyniki analizatorów) \cite{Jaoua2025LLMStaticAnalyzersCR} albo jako dostarczanie kuratorowanej "mapy wiedzy", 
która działa bardziej jak element rozumowania symbolicznego \cite{Icoz2025SymbolicReasoning}. W skrócie: PEFT pomaga, 
tylko nie rozwiązuje wszystkiego sam, a w części przypadków sensownie jest go 
łączyć z retrieval lub dodatkowymi źródłami sygnału.
